\section*{СПИСОК ИСПОЛЬЗОВАННОЙ ЛИТЕРАТУРЫ}
\addcontentsline{toc}{section}{СПИСОК ИСПОЛЬЗОВАННОЙ ЛИТЕРАТУРЫ}

\begin{enumerate}
    \item Агаджанян Н.А. Экология человека: Учебник для вузов / Н.А. Агаджанян, Ю.П. Гичев, В.И. Торшин. – М.: КРУК, 2017. – 320 с.
    \item Рахманин Ю.А. Качество питьевой воды и здоровье населения мегаполисов // Гигиена и санитария. – 2018. – № 12. – С. 5–10.
    \item Wright S.L., Kelly F.J. Plastic and Human Health: A Micro Issue? // Environmental Science \& Technology. – 2017. – Vol. 51, No. 12. – P. 6634-6647.
    \item Chang A.M. et al. Evening use of light-emitting eReaders negatively affects sleep, circadian timing, and next-morning alertness // Proceedings of the National Academy of Sciences. – 2015. – Vol. 112, No. 4. – P. 1232-1237.
    \item Документация Telegram Bot API. [Электронный ресурс]. URL: https://core.telegram.org/bots/api (дата обращения: 10.05.2025).
    \item Официальная документация библиотеки pyTelegramBotAPI (Telebot). [Электронный ресурс]. URL: https://github.com/eternnoir/pyTelegramBotAPI (дата обращения: 11.05.2025).
    \item Руководство по HTML и CSS от Mozilla Developer Network (MDN). [Электронный ресурс]. URL: https://developer.mozilla.org/ru/ (дата обращения: 12.05.2025).
    \item Марк Лутц. Изучаем Python. 5-е издание. – СПб.: Символ-Плюс, 2019. – 832 с.
    \item Фримен Э., Робсон Э. Изучаем HTML, XHTML и CSS. – СПб.: Питер, 2018. – 720 с.
    \item Кузнецов М.В. Git. Современная система контроля версий. – М.: ДМК Пресс, 2021. – 250 с.
    \item ВОЗ: Информационный бюллетень о микропластике в питьевой воде. [Электронный ресурс]. URL: https://www.who.int/water\_sanitation\_health/water-quality/guidelines/microplastics-in-drinking-water/ru/ (дата обращения: 15.05.2025).
    \item Паттерны проектирования: Конечный автомат (Finite State Machine). [Электронный ресурс]. URL: https://refactoring.guru/ru/design-patterns/state (дата обращения: 16.05.2025).
\end{enumerate}
